%======================================================================
\section{The fifth layer: the switch calculus transfers}\label{sec:fifth}

The remark closing Section~\ref{sec:switch} promised that nothing in the
switch calculus is specific to \emph{two} added rows.  This section makes
that promise concrete for the fifth layer: we derive the triple
inclusion--exclusion identity, transfer the switch decomposition and the
bracket calculus, and obtain
\begin{equation}\label{eq:r5thm}
r_5(n)=r_3(n)\bigl(1+O(n^{-1})\bigr),
\end{equation}
where $r_5$ is the fifth-layer effect
$\Delta_\lambda\log\bigl(N^{(5)}/N^{(4)}\bigr)$ between types $(2,n{-}2)$
and $(n)$.  Exact computations (independent of the calculus) confirm the
result at $n=8,9$ with striking precision.  Everything is written at the
same working-note rigor as Section~\ref{sec:switch}: exact identities and
finite verifications are labelled as such, and the two genuinely new
bookkeeping steps (the spectator factors in the strata analysis) are
flagged explicitly.

\subsection{A triple inclusion--exclusion identity}

$N^{(5)}(\sigma)$ is the number of ordered triples $(\tau,\eta,\theta)$
of $B_2$-avoiding permutations that are pairwise cell-disjoint.
Expanding each of the three pairwise-disjointness indicators by
inclusion--exclusion over subsets of that pair's coincidence set --- the
three expansions are independent --- gives the analogue of
Proposition~\ref{prop:pairIE}.

\begin{proposition}[triple inclusion--exclusion]\label{prop:tripleIE}
Let $M_1,M_2,M_3$ range over matchings of the complement of $B_2$
(the coincidence sets of $(\tau,\eta)$, $(\tau,\theta)$,
$(\eta,\theta)$ respectively), subject to \emph{compatibility}: each
pairwise union $M_i\cup M_j$ is again a matching (shared cells are
allowed; two distinct cells of $M_i$, $M_j$ may not share a line).
Then
\begin{equation}\label{eq:tripleIE}
N^{(5)}(\sigma)=\sum_{(M_1,M_2,M_3)}(-1)^{|M_1|+|M_2|+|M_3|}
R(M_1{\cup}M_2)\,R(M_1{\cup}M_3)\,R(M_2{\cup}M_3),
\end{equation}
with $R(\cdot)$ the punctured counts of \eqref{eq:punctured}.
Equivalently, in resummed form: compatibility of the triple is
equivalent to $U=M_1\cup M_2\cup M_3$ being a matching, and collecting
the label multiplicities cell by cell,
\begin{equation}\label{eq:tripleIEU}
N^{(5)}(\sigma)=\sum_{U}\ \sum_{U=U_{ab}\sqcup U_{ac}\sqcup U_{bc}\sqcup U_{abc}}
(-1)^{|U_{ab}|+|U_{ac}|+|U_{bc}|}\,2^{|U_{abc}|}\,
R(A)\,R(B)\,R(C),
\end{equation}
where $A=U_{ab}\cup U_{ac}\cup U_{abc}$,
$B=U_{ab}\cup U_{bc}\cup U_{abc}$, $C=U_{ac}\cup U_{bc}\cup U_{abc}$
are the constraint sets of $\tau,\eta,\theta$: \emph{every coincidence
cell lies in at least two of the three constraint sets}.
\end{proposition}

\begin{proof}
For fixed $\tau,\eta,\theta$ avoiding $B_2$, the product of the three
indicators $\prod_{\text{pairs}}\mathbf 1[\text{disjoint}]$ expands as
$\sum_{M_1\subseteq\tau\cap\eta}(-1)^{|M_1|}\cdots$, independently in
the three pairs; exchanging sums gives \eqref{eq:tripleIE}, the count of
triples with $\tau\supseteq M_1\cup M_2$ etc.\ factorizing since the
three rows are constrained separately (and vanishing unless each union
is a matching).  For \eqref{eq:tripleIEU}: if two cells of the $M_i$
share a line, compatibility of the relevant pair forces them equal, so
pairwise compatibility is the same as $U$ being a matching; each
$u\in U$ carries the nonempty label set $T(u)\subseteq\{1,2,3\}$ of the
$M_i$ containing it, the constraint-set membership of $u$ depends only
on $T(u)$ ($|T|=1$ gives exactly two memberships, $|T|\ge2$ gives all
three), and summing $(-1)^{|T|}$ over the label sets with a given
membership pattern gives weight $-1$ for each two-membership pattern
and $+3-1=+2$ for the full pattern.
(Both forms verified against brute-force enumeration of all triples at
$n=5$, and \eqref{eq:tripleIEU} additionally at $n=6,7$, both cycle
types: $N^{(5)}$ values $36$; $4032$, $4224$; $2185152$, $2198160$.
\texttt{triple\_ie.py}.)
\end{proof}

\begin{remark}[the truncated series is slowly convergent]
Truncating \eqref{eq:tripleIEU} at total order
$q=|M_1|+|M_2|+|M_3|\le Q$ gives an alternating series whose natural
parameter is the \emph{total} coincidence intensity of the three pairs,
$3\mu\approx6$ --- three times that of the pair series of
Section~\ref{sec:pair}.  Convergence in $Q$ is correspondingly slower:
at $n=8$ the $Q\le3$ truncation underestimates the fifth-layer effect
by $24\%$ (exact anchor below), at $n=9$ by $6\%$; the truncation error
shrinks rapidly with $n$ ($Q=3$ values $0.922$, $0.942$ at $n=10,11$ in
the $r_5\cdot(n)_4$ normalisation of the table below, each presumably a
few percent low).  We therefore use the truncated series only as a
consistency check at larger $n$ and rely on exact enumeration
(\texttt{n5\_direct2.c}) at $n\le9$; the switch calculus below is what
replaces series estimation as the source of the theorem.
(\texttt{layer5\_production.py}.)
\end{remark}

\subsection{The switch decomposition for triples}

Adopt the setting of Section~\ref{sec:switch}: adjacent representatives
$\sigma,\sigma'$, common board $B_0$ (two paths), pins
$d_1=(2,3)$, $d_2=(n,1)$, $e_1=(2,1)$, $e_2=(n,3)$, and $A_0$ the set of
$B_0$-avoiding permutations.  Write
\[
N^{(5)}_0=\#\{(\tau,\eta,\theta)\in A_0^3\ \text{pairwise disjoint}\},
\qquad
G^{(5)}(z)=\#\{\text{such triples with }z\in\tau\},
\]
and analogously $G_2^{(5)}(z,z')$ ($z,z'\in\tau$) and $H^{(5)}(z,z')$
($z\in\tau$, $z'\in\eta$).

\begin{lemma}[switch decomposition, fifth layer; exact]\label{lem:switchIE5}
With $\Delta N^{(5)}=N^{(5)}(\sigma')-N^{(5)}(\sigma)$,
\begin{align}
\Delta N^{(5)}
&=3\,\mathrm{Br}^{(5)}_G
+3\bigl[G_2^{(5)}(e_1,e_2)-G_2^{(5)}(d_1,d_2)\bigr]
+6\bigl[H^{(5)}(e_1,e_2)-H^{(5)}(d_1,d_2)\bigr],
\label{eq:switchIE5}\\
\mathrm{Br}^{(5)}_G&:=G^{(5)}(d_1)+G^{(5)}(d_2)-G^{(5)}(e_1)-G^{(5)}(e_2).
\notag
\end{align}
\end{lemma}

\begin{proof}
As for Lemma~\ref{lem:switchIE}: inclusion--exclusion over the events
$d_i\in\tau\cup\eta\cup\theta$, with $U(z)=3G^{(5)}(z)$ (choice of the
row through $z$) and $U(z,z')=3G_2^{(5)}(z,z')+6H^{(5)}(z,z')$ (one row
through both, or an ordered pair of distinct rows).  Verified by
complete enumeration at $n=6$
($\Delta N^{(5)}=192=288-24-72$) and $n=7$
($13008=15384+2496-4872$); and at $n=8$ the pieces computed on $B_0$
(\texttt{anatomy5.c}: $3\mathrm{Br}^{(5)}_G=2649312$,
$3\,\Delta G_2^{(5)}=159840$, $6\,\Delta H^{(5)}=-596544$) sum to
$2212608$, exactly the value of $\Delta N^{(5)}$ from the independent
direct enumeration of $N^{(5)}$ (\texttt{n5\_direct2.c}).
\end{proof}

\begin{remark}[$\Psi$ acts diagonally]
The involution $\Psi$ of Remark~\ref{rem:psi} preserves $A_0$ and
disjointness, hence acts on triples coordinatewise; consequently
$G^{(5)}(d_1)=G^{(5)}(d_2)$ \emph{exactly} (confirmed at $n=7,8,9$:
e.g.\ $G^{(5)}(d_1)=G^{(5)}(d_2)=496941376$ at $n=8$).
\end{remark}

\subsection{The stratified bracket; spectator factors}

Expand $G^{(5)}(z)$ by Proposition~\ref{prop:tripleIE} on the board
$B_0$ with the pin $z$ carried by $\tau$'s factor.  In the notation of
\eqref{eq:tripleIEU} ($A$ = $\tau$'s constraint set, $B,C$ the other
two),
\begin{equation}\label{eq:BrG5strata}
\mathrm{Br}^{(5)}_G=\sum_{A}\,B_A\,w(A),
\qquad
w(A):=\!\!\sum_{\substack{(M_1,M_2,M_3)\\ M_1\cup M_2=A}}\!\!
(-1)^{|M_1|+|M_2|+|M_3|}R_0(B)\,R_0(C),
\end{equation}
where $B_A=\sum_{z\in\{d_1,d_2\}}R_0^\dagger(A\cup z)
-\sum_{z\in\{e_1,e_2\}}R_0^\dagger(A\cup z)$ is \emph{the same}
four-term bracket as in \eqref{eq:strata}, now indexed by $A$.  Two
exact identities anchor \eqref{eq:BrG5strata}:
\begin{itemize}[leftmargin=2em]
\item[(i)] $\sum_A R_0(A)\,w(A)=N^{(5)}_0$ (the regrouped
\eqref{eq:tripleIEU} on $B_0$);
\item[(ii)] $w(\emptyset)=\sum_{M_3}(-1)^{|M_3|}R_0(M_3)^2=N^{(4)}_0$
(the pair inclusion--exclusion on $B_0$), so the $A=\emptyset$ stratum
equals $M_{n-4}\,N^{(4)}_0$ \emph{exactly} by
Proposition~\ref{prop:S0}.
\end{itemize}
Since $B_A$ is unchanged, the entire bracket toolkit of
Section~\ref{sec:switch} --- Lemma~\ref{lem:wronskian},
Lemma~\ref{lem:valuation}, Lemma~\ref{lem:master},
Proposition~\ref{prop:generic} --- applies verbatim in the $A$-slot.
What is new is the weight: $R_0(M)$ is replaced by the spectator weight
$w(A)$.  The transfer rests on two observations.

\emph{First observation.}  Every cell of $A$ lies in $B$ or $C$
(cells of $M_1$ lie in $B$, cells of $M_2$ in $C$).  Hence a line-swap
pairing $\varphi$ moving a cell of $A$ (as in Lemma~\ref{lem:pairing}
or Lemma~\ref{lem:colexchange}) moves the \emph{same} cell in exactly
one spectator factor when the cell has a single label, and in both when
it is shared; the third factor is untouched in the first case.  The
pinned-factor equalities of the pairing lemmas --- which involve only
$B_A$, through the deleted row- and column-\emph{sets} of the pinned
counts --- are therefore untouched, while the free-factor difference
becomes, by the product rule,
\[
R_0(B)R_0(C)-R_0(\varphi B)R_0(\varphi C)
=\bigl[R_0(B)-R_0(\varphi B)\bigr]R_0(C)
+R_0(\varphi B)\bigl[R_0(C)-R_0(\varphi C)\bigr],
\]
each bracket a single line-swap difference controlled by
Lemma~\ref{lem:wronskian} exactly as at the fourth layer, i.e.\ a
relative $O(n^{-3})$ per factor.

\emph{Second observation (spectator stability).}  The pairing must also
re-index the spectator sum over $M_3$: an $M_3$ compatible with
$(M_1,M_2)$ maps to one compatible with $\varphi(M_1,M_2)$ by the
identity map except on the $O(1/n)$-measure set of $M_3$ meeting the
two swapped lines, which is absorbed by (the $B_0$-analogue of)
Lemma~\ref{lem:linemeasure} applied inside $w$.  Thus
$|w(A)-w(\varphi A)|=O(n^{-1})\,\overline w(A)$ with
$\overline w$ the unsigned weight, and the paired-class cancellations
of Lemmas~\ref{lem:pairing}, \ref{lem:Hdiff}, \ref{lem:colexchange}
and~\ref{lem:boundary} survive with the same $O(n^{-1})$ headroom.
Finally, Lemma~\ref{lem:uniform} and Corollary~\ref{cor:moments}
apply to $R_0(B)$, $R_0(C)$ individually (any max-degree-2 board), so
the unsigned sums $\sum_A|B_A|\,\overline w(A)$ obey the same stratum
bounds as at the fourth layer with $N_0^2$ replaced by an
$N^{(5)}_0$-scale; the $G_2$- and $H$-differences of
\eqref{eq:switchIE5} are handled by the same pairings with one or two
pins carried by spectator slots.

These two observations are the only new content; we record the outcome.

\begin{theorem}[fifth layer]\label{thm:layer5}
For $n\ge n_0$,
\[
\log\frac{N^{(5)}_n((2,n{-}2))}{N^{(5)}_n((n))}
=3\,\frac{M_{n-4}}{M_n}+O(n^{-5}),
\qquad\text{hence}\qquad
r_5(n)=r_4(n)\bigl(1+O(n^{-1})\bigr)=r_3(n)\bigl(1+O(n^{-1})\bigr).
\]
\end{theorem}

\begin{proof}
Run the proof of Theorem~\ref{thm:layer4} on \eqref{eq:switchIE5} and
\eqref{eq:BrG5strata}, with Proposition~\ref{prop:generic} in the
$A$-slot, weights $w(A)$ in place of $R_0(M)$, and the two observations
above supplying the paired-class and unsigned-sum bounds; the generic
strata resum, via identity (i), to
$\mathrm{Br}^{(5)}_G=(M_{n-4}/M_n)\,N^{(5)}_0\bigl(1+O(1/n)\bigr)$, and
the pin-pin differences of \eqref{eq:switchIE5} are
$O\bigl(n^{-1}(M_{n-4}/M_n)N^{(5)}_0\bigr)$.  Then
$N^{(5)}_0=N^{(5)}((n))(1+O(1/n))$ as before, giving
$\Delta\log N^{(5)}=3M_{n-4}/M_n+O(n^{-5})$, and
$r_5-r_3=\Delta\log N^{(5)}-\Delta\log N^{(4)}-\Delta\log N^{(3)}
=O(n^{-5})$.
\end{proof}

\subsection{Exact fifth-layer data}

Direct enumeration (64/128-bit cell masks; \texttt{n5\_direct.c},
\texttt{n5\_direct2.c}) gives $N^{(5)}$ exactly at $n\le9$; the $N^{(3)}$
and $N^{(4)}$ values reproduce the known exact tables, certifying the
enumeration.

\begin{center}
\begin{tabular}{r|rr|rrr|rr}
$n$ & $N^{(5)}((n))$ & $N^{(5)}((2,n{-}2))$ &
$r_3(n)_4$ & $r_4(n)_4$ & $r_5(n)_4$ &
$\frac{\Delta\log N^{(4)}}{2M_{n-4}/M_n}$ &
$\frac{\Delta\log N^{(5)}}{3M_{n-4}/M_n}$\\\hline
$6$ & $4032$ & $4224$ & $0$ & $2.454$ & $14.29$ & --- & ---\\
$7$ & $2185152$ & $2198160$ & $1.450$ & $2.560$ & $0.976$ & $1.382$ & $1.146$\\
$8$ & $1539331584$ & $1541544192$ & $0.709$ & $0.842$ & $0.862$ & $1.094$ & $1.134$\\
$9$ & $1417404643008$ & $1418731458816$ & $0.906$ & $0.948$ & $0.976$ & $1.023$ & $1.041$\\
\end{tabular}
\end{center}

\noindent Here $r_k(n)_4$ abbreviates $r_k\cdot n(n{-}1)(n{-}2)(n{-}3)$,
and $r_k=\Delta_\lambda\log(N^{(k)}/N^{(k-1)})$.  At $n=8,9$:
$r_5/r_4=1.023$, $1.030$ --- the fifth-layer effect tracks the fourth
to $3\%$ already at $n=9$ (small $n=6,7$ are off-trend, as for every
other layer statistic).  On the anatomy side,
$\mathrm{Br}^{(5)}_G=5128$, $883104$ at $n=7,8$ with
$\mathrm{Br}^{(5)}_G N_0/(M_{n-4}N^{(5)}_0)=0.612$, $0.696$,
consistent with a $1-c/n$ approach with $c\approx2.4$ ---
about twice the fourth-layer constant $c\approx1.17$, as the two
spectator rows suggest.

\begin{remark}[general $k$; the profile at the origin]\label{rem:layerk}
By induction on $k$, the identical transfer (a $(k{-}2)$-tuple
inclusion--exclusion; one more spectator factor per layer) gives, for
each \emph{fixed} $k$,
\[
\Delta_\lambda\log N^{(k)}=(k-2)\,\frac{M_{n-4}}{M_n}+O_k(n^{-5}),
\qquad
r_k(n)=r_3(n)\bigl(1+O_k(1/n)\bigr),
\]
with constants growing with $k$ (each spectator contributes its own
$O(1/n)$ bookkeeping).  In the language of the layer profile
$\gamma$ of Section~\ref{sec:layers}, $\gamma$ is flat at the origin:
every fixed layer contributes the same $(1+o(1))\,M_{n-4}/M_n$ per
$2$-cycle.  The uniformity in $k$ needed for the bulk of the profile
(layers $k\asymp tn$) remains the province of the cluster expansion
(TODO~4): there the boards are no longer two paths and the calculus
here does not apply.
\end{remark}
